% Generated draft: requires mathematical and source review.
\documentclass[11pt]{article}
\usepackage[margin=25mm]{geometry}
\usepackage{fontspec}
\setmainfont{Latin Modern Roman}
\usepackage{amsmath,amssymb,mathtools,unicode-math}
\setmathfont{Latin Modern Math}
\usepackage{enumitem,needspace,xurl,hyperref}
\setlength{\parindent}{0pt}
\setlength{\parskip}{5pt}
\setlength{\emergencystretch}{3em}
\newcommand{\N}{\mathbb{N}}
\newcommand{\Z}{\mathbb{Z}}
\newcommand{\Q}{\mathbb{Q}}
\newcommand{\R}{\mathbb{R}}
\newcommand{\C}{\mathbb{C}}
\newcommand{\F}{\mathbb{F}}
\newcommand{\A}{\mathbb{A}}
\newcommand{\PP}{\mathbb{P}}
\newcommand{\E}{\mathbb{E}}
\newcommand{\eps}{\varepsilon}
\newcommand{\dd}{\,\mathrm{d}}
\begin{document}
\section*{UnsolvedMath ID: 800941}
\textbf{Source problem code:} catalogue entry
\subsection{Definitions and mathematical statement}
Under the conventional Riemannian interpretation, let \((M,g)\) be a smooth Riemannian manifold, let \(SM=\{v\in TM:g(v,v)=1\}\), and write \(\varphi_g^t\) for its unit-speed geodesic flow. Negative curvature means \(K_g(\sigma)<0\) for every two-dimensional tangent plane \(\sigma\). The supplied record does not explicitly specify this interpretation, the dimension, compactness, completeness, or regularity assumptions on the metric.

At a Lyapunov-regular vector \(v\in SM\), consider the finite Lyapunov exponents of \(D\varphi_g^t\), computed using the Sasaki norm induced by \(g\). When the positive exponents form a nonempty set, define
\[
\lambda_{\max}^{+}(v)=\max\{\lambda:\lambda>0\text{ is a Lyapunov exponent at }v\},\qquad
\lambda_{\min}^{+}(v)=\min\{\lambda:\lambda>0\text{ is a Lyapunov exponent at }v\}.
\]
The question attributed in the supplied record to Burns and Matveev is whether there exists a negatively curved metric for which
\[
\frac{\lambda_{\max}^{+}(v)}{\lambda_{\min}^{+}(v)}>2
\]
for typical geodesics, rather than only for exceptional closed geodesics.

The word ``typical'' is unresolved in the supplied record. If it means almost everywhere for a specified invariant probability measure \(\mu_g\) on \(SM\), the corresponding requirement is
\[
\mu_g\!\left(\left\{v\in SM:\begin{array}{l}
\varphi_g^t(v)\text{ is defined for all }t\in\R,\\
v\text{ is Lyapunov regular with finite exponents},\\
0<\lambda_{\min}^{+}(v)\leq\lambda_{\max}^{+}(v),\\
\lambda_{\max}^{+}(v)>2\lambda_{\min}^{+}(v)
\end{array}\right\}\right)=1.
\]
This is a conditional measure-theoretic formulation, not a choice of the intended measure. A fully specified existence statement requires both the intended class of manifolds and metrics and the meaning of ``typical''; neither is supplied. No compactness, completeness, invariant measure, or almost-everywhere existence theorem for these exponents is assumed here.
\subsection{Short English statement}
Does some negatively curved metric have largest positive geodesic-flow Lyapunov exponent strictly more than twice its smallest positive exponent along typical geodesics, rather than merely along exceptional closed geodesics? The supplied statement leaves the ambient manifold class and the meaning of ``typical'' unspecified.
\subsection{Sources}
\begin{itemize}
\item Supplied local UnsolvedMath registry record: permanent ID \(800941\), exact source code \texttt{catalogue entry}, title \emph{Open Problems and Questions about Geodesics --- Question 21}. The registry attributes the question to Burns and Matveev, gives proposed year 2019, and labels it open. That label is supplied metadata, not an independently verified mathematical or literature assessment.
\item Supplied source attribution: Burns, Matveev, \emph{Open Problems and Questions about Geodesics} (2019), Question 21. Supplied URL: \url{https://arxiv.org/abs/1308.5417}. The background records access on 2026-10-16 and extraction by \texttt{arxiv\_latex}; these are supplied provenance claims, not independent source verification.
\end{itemize}
\end{document}
